\documentclass[12pt]{article}
\usepackage[utf8]{inputenc}
\usepackage[T1]{fontenc}
\usepackage{hyperref}
\title{Yitha Yitha Language Resources for K-12 Curriculum Development}
\author{Compiled by Claude Sonnet 5.5 (claude-sonnet-5-5) for Andr\'as Kornai}
\date{}
\begin{document}
\maketitle

\section{Introduction}
Yitha Yitha (ISO 639-3 xth; also Eetha-eetha, Ita-ita, Yetho, Yit-tha) is a dormant, unclassified Australian language. Wikipedia places it in southern New South Wales, while a mirror of the article gives South Australia, so the sources conflict. Glottolog notes that Yitha Yitha and Dadi Dadi are close enough to treat as one language, Lower Darling--Lower Lachlan. Almost nothing beyond a nineteenth-century word list was found.

\section{Primary Educational Resources}
None found. The language is not known to be taught in schools.

\section{Linguistic and Reference Materials}
\begin{itemize}
\item Wikipedia, ``Yitha Yitha language''. \url{https://en.wikipedia.org/wiki/Yitha_Yitha_language}
\item Glottolog reference: Macdonald 1886, ``Yit-tha'', in Curr, \textit{The Australian Race}, vol. II, pp. 288--289, the main historical word list. \url{https://glottolog.org/resource/reference/id/59575}
\item Glottolog language entry. \url{https://glottolog.org/resource/languoid/iso/xth}
\item Endangered Languages Project entry. \url{https://endangeredlanguages.com/lang/6747}
\item University of Queensland manuscripts record, ``Yitha Yitha language D7''. \url{https://manuscripts.library.uq.edu.au/index.php/yitha-yitha-language-d7}
\end{itemize}

\section{Religious and Community Resources}
None found.

\section{Digital and Audio Resources}
No audio or software was found. The State Library of Queensland word lists did not appear to include Yitha Yitha, so they are not listed as a resource.

\section{Curriculum Development Context}
Because the language is dormant and documented only by a short historical word list, a curriculum in the language is not realistic from current sources. Suitable work would be awareness and heritage content on the Lower Darling and Lachlan region, done with the relevant Aboriginal communities and language bodies. The Macdonald 1886 list and the UQ manuscript record are the first sources to examine, and any revival should be led by community members.
\end{document}